\documentclass[11pt,letterpaper]{article}

\usepackage[margin=0.65in]{geometry}
\usepackage{enumitem}
\usepackage{hyperref}
\usepackage{titlesec}
\usepackage{tabularx}
\usepackage{xcolor}
\usepackage[T1]{fontenc}
\usepackage[utf8]{inputenc}
\usepackage{charter}

\hypersetup{
    colorlinks=true,
    urlcolor=blue,
    linkcolor=black
}

\pagestyle{empty}
\setlength{\parindent}{0pt}
\setlength{\parskip}{0pt}
\setlist[itemize]{leftmargin=1.2em, itemsep=1pt, topsep=2pt, parsep=0pt, partopsep=0pt}

% section style
\titleformat{\section}
    {\bfseries\large}
    {}
    {0em}
    {}
    [\titlerule]

\titlespacing*{\section}{0pt}{8pt}{4pt}

% custom entry macro
\newcommand{\resumeSubheading}[4]{
    \begin{tabularx}{\textwidth}{@{}Xr@{}}
        \textbf{#1} & \textbf{#2} \\
        #3 & #4 \\
    \end{tabularx}
}

\newcommand{\resumeItem}[1]{
    \item #1
}

\begin{document}

%==================== HEADER ====================
\begin{center}
    {\LARGE \textbf{Ziyu (Rick) Xu}} \\[2pt]
    {\small Robotics Institute, Carnegie Mellon University} \\[1pt]
    {\small \textbf{Research Interests:} Robot Foundation Models, Real-World Policy Adaptation, and Learning-Augmented Control} \\[1pt]
    {\small
    +1 7348005258
    \,|\, 
    \href{mailto:rickx@andrew.cmu.edu}{rickx@andrew.cmu.edu}
    } \\[1pt]
\end{center}

%==================== EDUCATION ====================
\section*{Education}

\resumeSubheading
{Carnegie Mellon University (CMU)}
{Aug. 2026}
{\hspace*{1.4em}Master of Science in Robotics (MSR)}
{Pittsburgh, PA}
\resumeSubheading
{University of Michigan, Ann Arbor (UM)}
{Oct. 2024 -- May 2026}
{\hspace*{1.4em}B.S.E. in Mechanical Engineering (Robotics Track); Minor in Computer Science}
{GPA: 3.85/4.0 (Top 5\%)}
\begin{itemize}[leftmargin=2.4em, itemsep=0.5pt, topsep=0pt]
    \item \textbf{Core Coursework:} Robot Learning, RL, Embedded Systems, Optimal Control, Convex Optimization
\end{itemize}

\resumeSubheading
{Shanghai Jiao Tong University (SJTU)}
{Sep. 2022 -- Jul. 2026}
{\hspace*{1.4em}B.S.E. in Electrical and Computer Engineering}
{GPA: 3.56/4.0 (Top 20\%)}
\begin{itemize}[leftmargin=2.4em, itemsep=0.5pt, topsep=0pt]
    \item \textbf{Core Coursework:} Algorithms, Mathematical Foundations, Probability Theory, Signals and Systems
\end{itemize}
%==================== RESEARCH ====================
\section*{Research Experiences and Internships}

\resumeSubheading
{ Unitree Robotics, R\&D Department, Embodied AI Team}
{May 2026 -- Present}
{\hspace*{1.4em}Research Intern, Humanoid Robot Foundation Models}
{Hangzhou, China}
\begin{itemize}[leftmargin=2.4em, itemsep=0.5pt, topsep=0pt]
    \small
    \resumeItem{\textbf{\textit{Closed-Loop Real-Robot Infrastructure for Humanoid Foundation-Model Policies:}} Built and validated an end-to-end policy infrastructure on Unitree H2/G1 humanoids, integrating teleoperation data collection, VLA/WAM and whole-body policy adapters, rollout logging and evaluation, controller interfaces, and policy fine-tuning. Enabled launch with one command and closed-loop rollouts across platforms, allowing policies and controllers to be integrated without modifying platform-specific launch workflows.}
    \resumeItem{\textbf{\textit{Embodiment VLA Adaptation \& Post-Training:}} Conducted a structured survey of representative robot VLA, WAM, and physical-AI policies, especially work from NVIDIA GEAR Lab. Deployed and fine-tuned Isaac GR00T/OpenPI VLA policies with SONIC whole-body control on Unitree humanoids using real-robot rollout traces. Currently developing a phase-aware adaptation method that aligns semantic task progress and execution timing across embodiments for policy transfer and rollout-driven post-training.}
    \resumeItem{\textbf{\textit{Reproducible Real-World VLA Benchmark for Humanoid Manipulation:}} Implemented an 8-task, 24-scene benchmark with green-screen background replacement to evaluate foundation models under varied visual conditions. Built a PICO-based scene-digitization and AR-guided reconstruction pipeline that captured shelf-object layouts as reusable digital scene templates, and registered target placements back into the physical workspace, achieving 2.5 cm median placement error.}
\end{itemize}

\resumeSubheading
{Shanghai Artificial Intelligence Lab \& IRMV, SJTU }
{Apr. 2026 -- Present}
{\hspace*{1.4em}Research Assistant, Directed by Prof. Junchi Yan and Prof. Hesheng Wang}
{Shanghai, China}
\begin{itemize}[leftmargin=2.4em, itemsep=0.5pt, topsep=0pt]
    \small
    \resumeItem{\textbf{\textit{Action-Step-Aware Dynamic Visual Sparsity for VLAs:}} Developing a fine-grained, action-step-aware visual-token pruning method that re-estimates token importance and retention budgets across action-generation steps rather than pruning once at prefill, preserving task-relevant visual evidence as manipulation context evolves, and preparing a first-author manuscript on efficient vision-language-action models.}
\end{itemize}

\resumeSubheading
{Computational Autonomy and Robotics Lab, UM }
{Jun. 2025 -- Apr. 2026}
{\hspace*{1.4em}Research Assistant, Directed by Prof. Maani Ghaffari}
{Ann Arbor, MI}
\begin{itemize}[leftmargin=2.4em, itemsep=0.5pt, topsep=0pt]
    \small
    \resumeItem{\textbf{\textit{Online Learning Lie-MPC for Robust ASV Control:}} Developed a Neural-Fly-inspired learning-augmented Lie-MPC framework for an autonomous surface vehicle, estimating residual wind/current disturbances online and injecting learned dynamics corrections into the controller for robust trajectory tracking under hydrodynamic uncertainty. Implemented and evaluated the framework on a real ASV, achieved a 20\% improvement in trajectory-tracking performance over the baseline MPC controller during field experiments.}
    \resumeItem{\textbf{\textit{Field Perception for Aquatic Vegetation Tracking:}} Built a water-scene perception pipeline for autonomous surface vehicles, combining DINO/SAM2-based segmentation with Grounded-DINO and YOLO fine-tuning to detect submerged vegetation and provide perception signals for navigation and autonomy experiments.}
    \resumeItem{\textbf{\textit{ASV Hardware and Sensor Integration for Closed-Loop Field Autonomy:}} Designed sensor mounts and integrated GPS, lidar, IMU, and camera systems on an autonomous surface vehicle. Aligned and calibrated sensors to support field data collection and control experiments.}
\end{itemize}

\resumeSubheading
{Intelligent Robotics and Autonomy Lab , UM}
{Oct. 2024 -- Apr. 2026}
{\hspace*{1.4em}Research Assistant, Directed by Prof. Vasileios Tzoumas}
{Ann Arbor, MI}
\begin{itemize}[leftmargin=2.4em, itemsep=0.5pt, topsep=0pt]
    \small
    \resumeItem{\textbf{\textit{Multi-Robot Active-SLAM Simulation \& ROS2 Integration:}} Rebuilt a Unity + AirSim multi-drone simulation testbed with redesigned vehicle dynamics and sensing modules, then migrated SlideSLAM communication logics from ROS1 to ROS2 by refactoring topic interfaces and workflows for bandwidth-limited active-SLAM evaluation.}
    \resumeItem{\textbf{\textit{Resource-Aware Multi-Robot Exploration Algorithms:}} Developed distributed exploration EXP3-style algorithms for bandwidth-limited robot teams, formulating communication-aware action selection with submodular objectives and bandit feedback in simulation to balance map coverage and communication cost.}
\end{itemize}

\resumeSubheading
{Surgical Computer Vision Detection, SJTU}
{Apr. 2024 -- Dec. 2024}
{Research Assistant, Directed by Prof. Yutong Ban}
{Shanghai, China}
\begin{itemize}[leftmargin=2.4em, itemsep=0.5pt, topsep=0pt]
    \small
    \resumeItem{\textbf{\textit{Foundation Segmentation Model Evaluation for Surgical Video:}} Evaluated Segment Anything Model variants on surgical video frames, analyzing domain shift and temporal mask stability around instruments, tissues, and tool-tissue interaction regions.}
\end{itemize}

%==================== PUBLICATIONS ====================
\section*{Publications and Manuscripts}

\begin{itemize}[leftmargin=1.4em, itemsep=3pt, topsep=2pt]
    \item \textbf{Ziyu Xu}, Haoran Zhang, and Sen Peng. \\
    ``PhaseVLA: Semantic Phase and Speed Alignment for Cross-Embodiment VLA Models'', \\
    \textit{Manuscript in preparation}, 2026.
    \item \textbf{Ziyu Xu}, Xiaohan Wang, Anning Hu, Hesheng Wang, and Junchi Yan. \\
    ``Action-Step-Aware Dynamic Visual Sparsity for Efficient Vision-Language-Action Models'', \\
    \textit{Manuscript in preparation}, 2026.
    \item Yinan Dong, \textbf{Ziyu Xu}, Tsimafei Lazouski, Sangli Teng, and Maani Ghaffari. \\
    ``Online Learning-Enhanced Lie Algebraic MPC for Robust Autonomous Surface Vehicle Control'', \\
    \textit{IEEE Robotics and Automation Letters (RA-L)}, revised manuscript under review, 2025. \\
    \textit{Contribution:} Led the design and development of the learning-augmented Lie-MPC framework; contributed to real-ASV experiments and evaluation.
    \item C. Yuan, J. Jiang, K. Yang, \textbf{Z. Xu}, \ldots, and Y. Ban. \\
    ``Systematic Evaluation and Guidelines for Segment Anything Model in Surgical Video Analysis'', \
    \textit{npj Digital Surgery}, accepted, 2024.
    
\end{itemize}

%==================== ADDITIONAL EXPERIENCE ====================
\section*{Additional Experience}

\begin{itemize}[leftmargin=1.4em, itemsep=3pt, topsep=2pt]
    \small
    \resumeItem{\textbf{RL-Control led Multi-Branch Cooling System for Data center, Grundfos Company:} Designed an RL control algorithm for a multi-branch cooling system with heterogeneous thermal loads, regulating individual water-pump speeds to minimize energy consumption by 30\% relative to the provided trim-response algorithm.}
    \resumeItem{\textbf{Low-Cost Automated 96-Well Plate Washer:} Designed, built, and delivered a fully automated 96-well microplate washer to the University of Michigan Medical School at 10\% of the cost of a comparable commercial system.}
    \resumeItem{\textbf{Data Engineering \& Web Systems, Library of Shanghai Jiao Tong University:} Led a team of 4 to build a Python dataset and internal web platform supporting 4,000+ SJTU faculty.}
    \resumeItem{\textbf{Quantitative Analysis, APEXUS Tech LLC:} Developed a hybrid statistical-learning and factor-mining model, integrating it into an execution framework with reproducible experiment logs and adaptive parameter tuning.}
\end{itemize}

%==================== AWARDS ====================
\section*{Awards}

\begin{itemize}
    \resumeItem{Gold Medal in 2023 University Physics Competition, by American Physical Society \& American Astronomical Society (top 1\%)}
    \resumeItem{Third place in 2023 RoboMaster Campus Competition (top 10\%)}
    \resumeItem{Dean's Honor List, University of Michigan College of Engineering (2024, 2025)}
    \resumeItem{M Prize in 2024 Mathematical Contest in Modeling (top 13\%)}
\end{itemize}

%==================== SKILLS ====================
\section*{Skills}

\begin{itemize}
\resumeItem{\textbf{Robot Learning \& Control:} VLA post-training, humanoid foundation models, whole-body policy, RL, real-robot policy deployment, teleoperation/data collection, MPC, sim-to-real}
\resumeItem{\textbf{Perception \& Multimodal ML:} PyTorch, OpenCV, SAM2, Grounded-DINO, YOLO, DINO features, visual token compression, spatial grounding, LLMs fine-tuning}
\resumeItem{\textbf{Systems \& Tools:} ROS/ROS2, Python/C++, Unity + AirSim, Docker, Linux, Git, sensor integration, CAD, FEA, Manufacturing, embedded system}
\end{itemize}

\end{document}
